\section{GPU 架构快速回顾}

在进入本节课的核心内容之前，先快速回顾上一节课的关键概念。

\subsection{硬件结构}

GPU 的核心结构包括：
\begin{itemize}
    \item \textbf{SM（Streaming Multiprocessor）}：GPU 的计算单元，A100 有 108 个 SM
    \item \textbf{内存层次}：
    \begin{itemize}
        \item DRAM（全局内存）：A100 有 80GB，大但慢
        \item L2 Cache：A100 有 40MB
        \item L1 Cache / Shared Memory：A100 每个 SM 有 192KB，小但快
    \end{itemize}
\end{itemize}

\subsection{执行模型}

\begin{figure}[H]
\centering
\begin{tikzpicture}[
    box/.style={draw, rounded corners, minimum width=2.5cm, minimum height=0.8cm, font=\small},
    arr/.style={->, >=stealth, thick}
]
\node[box, fill=blue!10] (grid) {Grid（网格）};
\node[box, fill=green!10, below left=1.2cm and 0.5cm of grid] (block0) {Block 0};
\node[box, fill=green!10, below=1.2cm of grid] (block1) {Block 1};
\node[box, fill=green!10, below right=1.2cm and 0.5cm of grid] (blockn) {Block N};
\node[box, fill=orange!10, below=1.0cm of block1] (threads) {Thread 0, 1, ..., 31 (Warp)};

\draw[arr] (grid) -- (block0);
\draw[arr] (grid) -- (block1);
\draw[arr] (grid) -- (blockn);
\draw[arr] (block1) -- (threads);

\node[right=0.3cm of blockn, font=\scriptsize, text width=3.5cm] {每个 Block 调度到一个 SM，Block 内线程共享 Shared Memory};
\end{tikzpicture}
\caption{GPU 执行模型层次：Grid $\rightarrow$ Block $\rightarrow$ Thread}
\end{figure}

\begin{itemize}
    \item \textbf{Thread（线程）}：执行 $f(i)$ 的最小单位
    \item \textbf{Thread Block}（又称 CTA）：调度到单个 SM 上的线程集合，Block 内线程共享 Shared Memory
    \item \textbf{Grid}：所有 Thread Block 的集合
\end{itemize}

\subsection{Thread Block 的意义}

为什么需要 Thread Block 而不是直接使用裸线程？
\begin{enumerate}
    \item Block 内线程拥有\textbf{共享内存}（Shared Memory），速度与 L1 Cache 相当
    \item Block 内线程可以\textbf{同步}（synchronize），但跨 Block 无法同步
    \item 需要线程间通信的操作（如矩阵乘法）应安排在同一个 Block 内
\end{enumerate}

\subsection{Wave Quantization 与 SM 利用率}

Thread Block 以``波（wave）''为单位调度到 SM 上。理想情况下，每一波应填满所有 SM。

\begin{warningbox}{Wave Quantization 陷阱}
如果 Block 数不能被 SM 数量整除，最后一波会有 SM 空闲。经验法则：Block 数应 $\geq 4 \times$ SM 数量，以确保高利用率。
\end{warningbox}

\subsection{算术强度}

\begin{importantbox}{算术强度 = FLOPs / 字节}
\begin{itemize}
    \item 高算术强度 $\Rightarrow$ \textbf{计算受限}（compute-bound）——这是好事
    \item 低算术强度 $\Rightarrow$ \textbf{内存受限}（memory-bound）——GPU 计算单元闲置
\end{itemize}
一般规则：矩阵乘法是 compute-bound，\textbf{其余几乎所有操作都是 memory-bound}。
\end{importantbox}

\subsection{本章小结}

GPU 执行模型的核心是：写一个函数 $f(i)$，GPU 对所有 $i = 0, \dots, N-1$ 并行执行。Block 提供了快速的共享内存和同步机制，但跨 Block 通信代价高昂。性能优化的关键在于最小化全局内存读写、最大化算术强度。

%% ========== Section 3: Benchmarking ==========

